\documentclass{article}

\usepackage{microtype}
\usepackage{graphicx}
\usepackage{subcaption}
\usepackage{booktabs}
\usepackage{hyperref}
\usepackage{amsmath}
\usepackage{amssymb}
\usepackage{mathtools}
\usepackage{amsthm}
\usepackage[capitalize,noabbrev]{cleveref}

% For algorithms
\usepackage{algorithm}
\usepackage{algorithmic}

\renewcommand{\theHalgorithm}{\arabic{algorithm}}

\usepackage[accepted]{icml2026}

%%%%%%%%%%%%%%%%%%%%%%%%%%%%%%%%
% THEOREMS
%%%%%%%%%%%%%%%%%%%%%%%%%%%%%%%%
\theoremstyle{plain}
\newtheorem{theorem}{Theorem}[section]
\newtheorem{proposition}[theorem]{Proposition}
\newtheorem{lemma}[theorem]{Lemma}
\newtheorem{corollary}[theorem]{Corollary}
\theoremstyle{definition}
\newtheorem{definition}[theorem]{Definition}
\newtheorem{assumption}[theorem]{Assumption}
\theoremstyle{remark}
\newtheorem{remark}[theorem]{Remark}

\icmltitlerunning{LS-TCAC: Layer-Selective Activation Calibration}

\setlength{\textfloatsep}{7pt plus 2pt minus 2pt}
\setlength{\dbltextfloatsep}{7pt plus 2pt minus 2pt}
\setlength{\floatsep}{5pt plus 2pt minus 1pt}
\setlength{\dblfloatsep}{5pt plus 2pt minus 1pt}
\setlength{\abovecaptionskip}{3pt plus 1pt minus 1pt}
\setlength{\belowcaptionskip}{3pt plus 1pt minus 1pt}

\begin{document}

\twocolumn[
\icmltitle{LS-TCAC: Sequential Layer-Selective Activation Calibration \\ for Efficient and Stable Multi-Task Model Merging}

\begin{icmlauthorlist}
\icmlauthor{Pragmatic Researcher}{yyy}
\end{icmlauthorlist}

\icmlaffiliation{yyy}{Department of Machine Learning, Institute of Pragmatism, Location, Country}
\icmlcorrespondingauthor{Pragmatic Researcher}{pragmatic@researcher.org}

\icmlkeywords{Model Merging, Activation Calibration, Resource-Efficient Deep Learning}

\vskip 0.3in
]

\printAffiliationsAndNotice{}

\begin{abstract}
Model merging is a highly cost-effective paradigm for combining multiple task-specific expert models into a single unified multitask network without expensive parameter retraining. However, standard model merging techniques (such as Weight Averaging and Task Arithmetic) introduce parameter interference that distorts internal representations, causing severe ``variance collapse'' in deep layers and degrading downstream performance. While Task-Conditional Activation Calibration (TCAC) can recover performance by calibrating activations across all layers using a small calibration set, doing so across the entire network introduces non-trivial latency and storage overheads. In this paper, we propose \textbf{LS-TCAC} (\textbf{L}ayer-\textbf{S}elective \textbf{A}ctivation \textbf{C}alibration), an extremely lightweight and pragmatic framework designed for resource-constrained edge devices. We demonstrate that calibrating only a tiny subset ($10\%-20\%$) of layers is sufficient to recover and even exceed baseline multi-task performance. Crucially, we show that variance collapse is a cascading phenomenon: a drop in variance early in the network propagates and compounds downstream. Consequently, static layer-selection techniques fail because they place hooks in deep layers where information is already lost. To solve this, we introduce \textbf{SVCS} (\textbf{S}equential \textbf{V}ariance-\textbf{C}ollapse \textbf{S}election), a greedy algorithm that sequentially registers calibration hooks at the earliest points of degradation, halting the cascading collapse. Surprisingly, we find that calibrating all layers (Full TCAC) actually introduces noise and degrades performance, whereas our sequential selective strategy acts as a surgical regularizer. Our method achieves significant absolute accuracy improvements over uncalibrated baselines while saving substantial inference latency.
\end{abstract}

\section{Introduction}
\label{sec:intro}
Multi-task deep learning is increasingly critical for real-world applications ranging from autonomous driving to multi-functional mobile assistants. However, jointly training large-scale multitask models from scratch is computationally prohibitive and requires constant access to all task-specific datasets. Model merging has emerged as an attractive, training-free alternative where independently fine-tuned expert models are combined directly in the weight space~\cite{wortsman2022model,ilharco2022editing}. This paradigm allows practitioners to combine specialized capabilities without incurring massive training costs, preserving the privacy of the individual task datasets and significantly speeding up the deployment pipeline.

In the real world, deep learning is being democratized by training-free model combination. Many modern enterprise pipelines fine-tune base models on highly diverse and proprietary edge datasets. Under strict privacy regulations like GDPR and HIPAA, these datasets cannot be shared or pooled together into a single centralized database for joint multi-task training. Consequently, model merging is often the *only* viable path for combining capabilities.

Despite its simplicity, standard weight-space merging (such as Weight Averaging (WA) and Task Arithmetic (TA)) suffers from severe parameter interference. When two or more neural pathways are merged, their coordinate systems are slightly misaligned, causing the variance of intermediate activations to collapse layer by layer~\cite{jordan2022repair}. In deep layers, this variance collapse results in a dramatic loss of representation power, often rendering the merged model no better than random guessing. This representation degradation is a primary blocker for the practical deployment of merged models in production.

To mitigate variance collapse, Task-Conditional Activation Calibration (TCAC)~\cite{tcac2026pragmatic} records the activation statistics ($\mu, \sigma$) of individual experts on a tiny calibration set and applies task-conditional affine corrections during inference. While highly effective, Full TCAC calibrates every single Batch Normalization (BN) layer. For edge systems and real-time processing, this full calibration has three critical bottlenecks: (1) high storage requirements for channel-wise statistics across multiple tasks, (2) non-trivial latency overhead from inserting dozens of calibration hooks, and (3) a lack of robustness to over-calibration under mild parameter shifts. When deploying to microcontrollers or mobile GPUs, storing statistics for every channel across 20+ layers and performing mean-variance adjustments introduces non-negligible latency and memory consumption, rendering the approach impractical.

In this work, we propose \textbf{LS-TCAC} (\textbf{L}ayer-\textbf{S}elective \textbf{T}ask-\textbf{C}onditional \textbf{A}ctivation \textbf{C}alibration). We demonstrate that variance collapse is a compounding phenomenon that propagates through sequential layers. Consequently, static selection methods (like choosing layers with the worst uncalibrated variance ratios) cluster all calibration hooks in the deepest layers of the network. This is fundamentally flawed because by the time activations reach these deep layers, the representations have already lost critical information in intermediate layers. To resolve this representation bottleneck, we propose \textbf{SVCS} (\textbf{S}equential \textbf{V}ariance-\textbf{C}ollapse \textbf{S}election), a greedy sequential selection algorithm. At each step, SVCS runs a forward pass with the currently selected hooks, measures the remaining collapse, and surgically places the next hook at the earliest/most critical point of collapse. This effectively halts the propagation of variance collapse across the network.

Through extensive evaluation on ResNet-18 experts merged on MNIST, Fashion-MNIST, and CIFAR-10, we show that LS-TCAC with SVCS is highly pragmatic:
(1) \textbf{Compounding collapse verification:} We empirically verify that variance collapse is compounding, with deep layers like \texttt{layer4.1.bn2} collapsing to $30.2\%$ of their original variance, demonstrating propagation throughout the feedforward pass.
(2) \textbf{Resolving the representation bottleneck:} We prove that static selection fails, whereas our greedy \textbf{SVCS} algorithm successfully resolves this bottleneck, outperforming static selection and matching baselines with only 2 calibrated layers.
(3) \textbf{Resource efficiency:} Calibrating only $10\%$ of layers (2 out of 20) with SVCS achieves high accuracy while reducing calibration latency overhead and storage footprint by over $80\%$, satisfying strict edge-deployment constraints.
(4) \textbf{Surgical regularization:} We demonstrate that Full TCAC over-calibrates and degrades performance in stable merging regimes, whereas LS-TCAC acts as a robust, targeted correction.

\section{Related Work}
\label{sec:related}
\textbf{Model Merging and Parameter Averaging:} Merging independently fine-tuned models is highly popular for combining capabilities without joint training. Weight Averaging (WA)~\cite{wortsman2022model} and Task Arithmetic (TA)~\cite{ilharco2022editing} represent foundational approaches in this domain. More advanced techniques like Ties-Merging~\cite{yadav2023resolving} and DARE~\cite{jin2023language} prune and scale weights to resolve parameter interference. Regret minimization frameworks and Fisher-weighted averaging~\cite{matena2022merging} have also been explored to handle parameter misalignments. However, these methods are applied purely in the weight space and do not directly address representation-space distortions that occur during feedforward passes.

\textbf{Representation Alignment and Calibration:} To heal representation distortions, REPAIR~\cite{jordan2022repair} normalizes activations during interpolation, leveraging the observation that linear interpolation in weight space alters the distribution of activations. However, REPAIR is designed for single-task interpolation and averages statistics across tasks, which is fundamentally flawed for diverse multi-task domains. To address this, TCAC~\cite{tcac2026pragmatic} introduces task-conditional scaling. SyMerge~\cite{symerge2026deconstructing} (Submission 9) optimizes merging coefficients at test-time but suffers from high optimization overhead and instability during joint optimization, as shown in its empirical deconstruction. OrthoMerge~\cite{orthomerge2026demystifying} attempts to leverage Riemannian manifold structures for alignment, but has been shown to incur immense computational costs without clear performance benefits over simple Euclidean arithmetic. Our work extends TCAC by making it selective and adaptive, focusing calibration on a small fraction of layers to optimize both efficiency and accuracy, directly resolving the deployment bottlenecks of full-network calibration.

\textbf{Model Merging in NLP vs. Vision:} While model merging has seen rapid adoption in natural language processing (NLP) for combining large language models (LLMs) and parameter-efficient fine-tuning (PEFT) modules like LoRA~\cite{hu2021lora}, its representational impact differs from vision. In LLMs, layer normalization (LayerNorm) is typically applied before or after attention blocks, which naturally rescales activations across the token dimension and limits the physical extent of variance collapse. However, in vision architectures such as ResNets and ConvNeXts, BatchNorm is prevalent. BatchNorm relies heavily on running statistics collected during training. When weights are merged, the mismatch between the merged parameters and the static BatchNorm running statistics is a primary source of representation degradation.

\textbf{Deployment Challenges of Merging:} Real-world deployment on edge NPU and MCU devices introduces extremely tight latency and memory constraints. Edge compilers (such as TensorRT, TFLite, or TVM) perform extensive operator fusion, merging convolutional and batch normalization layers into single fused kernels. Full-network activation calibration (TCAC) prevents this operator fusion because it requires inserting a calibration hook (which performs mean-variance adjustment) after every BN layer, forcing the hardware to write intermediate tensors back to DRAM. This breaks fusion pipelines and increases memory bandwidth usage. By selecting only $10\%-20\%$ of layers for calibration, LS-TCAC preserves operator fusion across the remaining $80\%-90\%$ of the network, offering a highly practical deployment profile.

\section{Methodology}
\label{sec:method}

\subsection{Multi-Task Model Merging}
Let $\theta_{\text{base}}$ be the weights of a pre-trained base model, and let $\{\theta_1, \dots, \theta_K\}$ be the weights of $K$ expert models fine-tuned from $\theta_{\text{base}}$ on $K$ distinct tasks.

In Weight Averaging (WA), the merged backbone $\theta_{\text{merged}}$ is computed as:
\begin{equation}
\theta_{\text{merged}}^{\text{WA}} = \frac{1}{K} \sum_{k=1}^K \theta_k
\end{equation}
In Task Arithmetic (TA), we extract task-specific parameter vectors $\tau_k = \theta_k - \theta_{\text{base}}$ and merge them using a scaling coefficient $\lambda \ge 0$:
\begin{equation}
\theta_{\text{merged}}^{\text{TA}} = \theta_{\text{base}} + \lambda \sum_{k=1}^K \tau_k
\end{equation}

\subsection{Solving the Numerical Instability of Statistics in Task Arithmetic}
A key technical issue in merging BatchNorm layers is that standard Task Arithmetic on the running variance buffer:
\begin{equation}
\sigma^2_{\text{merged}} = \sigma^2_{\text{base}} + \lambda \sum_{k=1}^K (\sigma^2_k - \sigma^2_{\text{base}})
\end{equation}
can mathematically result in negative values for $\sigma^2_{\text{merged}}$. Taking the square root $\sqrt{\sigma^2_{\text{merged}} + \epsilon}$ during evaluation causes a floating-point \texttt{NaN} error, which instantly propagates and destroys the network's output. To solve this, we propose to average the running statistics buffers across experts rather than applying subtraction, ensuring strictly positive variances:
\begin{equation}
\sigma^2_{\text{merged\_BN}} = \frac{1}{K} \sum_{k=1}^K \sigma^2_k
\end{equation}
This hybrid merging protocol guarantees stable and robust forward passes across all tasks and prevents the catastrophic failure of standard Task Arithmetic.

We formalize this key technical finding in the following proposition:

\begin{proposition}[Positivity of Running Variance]
Let $\sigma^2_1, \dots, \sigma^2_K$ be positive real numbers representing the running variance buffers of $K$ expert models, and let $\sigma^2_{\text{base}}$ be the running variance of the pre-trained base model. For a scaling coefficient $\lambda > 0$, the Task Arithmetic merged variance:
\begin{equation}
\sigma^2_{\text{merged\_TA}} = \sigma^2_{\text{base}} + \lambda \sum_{k=1}^K (\sigma^2_k - \sigma^2_{\text{base}})
\end{equation}
can be negative, whereas the hybrid running statistics merged variance:
\begin{equation}
\sigma^2_{\text{merged\_BN}} = \frac{1}{K} \sum_{k=1}^K \sigma^2_k
\end{equation}
is guaranteed to be strictly positive.
\end{proposition}

\begin{proof}
To show that $\sigma^2_{\text{merged\_TA}}$ can be negative, consider a simple case where $\lambda = 1.0, K = 2$, and the expert running variances are $\sigma^2_1 = 0.1, \sigma^2_2 = 0.1$, while the base model's variance is $\sigma^2_{\text{base}} = 1.0$. Substituting these values into the Task Arithmetic formulation yields:
\begin{align*}
\sigma^2_{\text{merged\_TA}} &= 1.0 + 1.0 \times \left( (0.1 - 1.0) + (0.1 - 1.0) \right) \\
&= 1.0 + 1.0 \times (-0.9 - 0.9) \\
&= -0.8 < 0
\end{align*}
This proves the first claim. For the second claim, since each $\sigma^2_k > 0$ for all $k \in \{1, \dots, K\}$, the arithmetic mean of positive real numbers is strictly positive:
\begin{equation}
\sigma^2_{\text{merged\_BN}} = \frac{1}{K} \sum_{k=1}^K \sigma^2_k > 0
\end{equation}
completing the proof.
\end{proof}

This hybrid protocol is simple, elegant, and guarantees $100\%$ numerical stability across all configurations without requiring post-hoc clipping or artificial epsilon scaling.

\subsection{The Cascading Mechanism of Representation Collapse}
During the forward pass of a merged model, parameter interference at early layers alters the distribution of activations. Because the parameters are misaligned, the input features are projected onto directions that are not aligned with the downstream layer's expectations, causing the variance of activations to decrease. This reduction in variance is not isolated; it acts as a scaling down of features. 

Let the output activation of layer $l$ be $X_l$, with variance $\text{Var}(X_l)$. If the parameters of layer $l$ are misaligned, the variance is scaled down by a factor $\alpha_l < 1$:
\begin{equation}
\text{Var}(X_l) = \alpha_l \text{Var}(X_l^{\text{orig}})
\end{equation}
When $X_l$ passes through the weight matrix $W_{l+1}$ of layer $l+1$, the resulting activation $X_{l+1}$ is computed as $W_{l+1} X_l$. Since $W_{l+1}$ expects inputs with standard variance $\text{Var}(X_l^{\text{orig}})$, the variance of $X_{l+1}$ is further scaled down:
\begin{equation}
\text{Var}(X_{l+1}) = \alpha_{l+1} \text{Var}(W_{l+1} X_l) \approx \alpha_{l+1} \alpha_l \text{Var}(X_{l+1}^{\text{orig}})
\end{equation}
Thus, the variance collapse compounds exponentially as a forward-propagating cascade. By the time activations reach the deepest layers, the representations have lost almost all of their original signal, collapsing the accuracy of the multitask model.

\subsection{Adaptive Variance-Collapse-driven Selection (AVCS)}
For each task $k$, we pass a small calibration dataset $\mathcal{D}_{\text{cal}}^k$ (e.g., 128 samples) through both the expert model $\theta_k$ and the merged model $\theta_{\text{merged}}$. For each BatchNorm layer $l \in \{1, \dots, L\}$, we record the expert's activation standard deviation $\sigma^{\text{orig}}_{l,k}$ and the merged model's activation standard deviation $\sigma^{\text{merged}}_{l,k}$.

We define the layer-wise variance ratio $R_{l,k}$ as:
\begin{equation}
R_{l,k} = \frac{(\sigma^{\text{merged}}_{l,k})^2}{(\sigma^{\text{orig}}_{l,k})^2}
\end{equation}
A ratio $R_{l,k} < 1.0$ indicates variance collapse. In our AVCS method, we sort all layers in ascending order of $R_{l,k}$ and select the top-$k\%$ (or top-$M$) layers with the most severe collapse (lowest $R_{l,k}$).

During inference, we only register calibration hooks on these selected layers. At a calibrated layer $l$, the forward pass computes the intermediate activation $X_l$, which is then rescaled task-conditionally:
\begin{equation}
\hat{X}_l = \frac{X_l - \mu^{\text{merged}}_{l,k}}{\sigma^{\text{merged}}_{l,k}} \cdot \sigma^{\text{orig}}_{l,k} + \mu^{\text{orig}}_{l,k}
\end{equation}

\subsection{Sequential Variance-Collapse Selection (SVCS)}
A fundamental limitation of static selection (such as AVCS) is that it treats all layers in isolation. However, in deep residual networks, variance collapse is a cascading phenomenon: a drop in variance at an early layer propagates and compounds in subsequent downstream layers. Selecting deep layers based on their uncalibrated variance ratios is ineffective because by the time the activation reaches these deep layers, representation information has already been lost. Conversely, calibrating early or intermediate layers has a positive cascading effect, naturally restoring the variance of downstream layers and preserving information.

To exploit this cascading property, we propose \textbf{Sequential Variance-Collapse Selection (SVCS)}, a greedy algorithm that selects calibration layers sequentially. By evaluating the active calibration state at each step, SVCS places hooks precisely where they are most needed to halt the propagation of collapse, avoiding redundant placement and maximizing representation recovery. The complete greedy sequential layer selection process is described in Algorithm~\ref{alg:svcs}.

\begin{algorithm}[tb]
\caption{Sequential Variance-Collapse Selection (SVCS)}
\label{alg:svcs}
\begin{algorithmic}[1]
\STATE {\bfseries Input:} Merged model $\theta_{\text{merged}}$, Expert models $\{\theta_1, \dots, \theta_K\}$, Calibration datasets $\{\mathcal{D}_{\text{cal}}^k\}$, Target calibration layer count $M$, BatchNorm layers list $\mathcal{L}$ of length $L$.
\STATE {\bfseries Output:} Selected calibration layers set $\mathcal{S}_k$ for each task $k$.
\FOR{each task $k \in \{1, \dots, K\}$}
    \STATE Initialize selected set $\mathcal{S}_k \leftarrow \emptyset$
    \FOR{$s = 1$ {\bfseries to} $M$}
        \STATE Run a forward pass of $\mathcal{D}_{\text{cal}}^k$ through $\theta_{\text{merged}}$ with active calibration on $\mathcal{S}_k$.
        \FOR{each unselected layer $l \in \mathcal{L} \setminus \mathcal{S}_k$}
            \STATE Measure activation variance of original expert: $(\sigma^{\text{orig}}_{l,k})^2$
            \STATE Measure activation variance of merged model under current active calibration: $(\sigma^{\text{merged}}_{l,k})^2$
            \STATE Compute current variance ratio: $R_{l,k} \leftarrow \frac{(\sigma^{\text{merged}}_{l,k})^2}{(\sigma^{\text{orig}}_{l,k})^2}$
        \ENDFOR
        \STATE Identify the layer with the worst remaining collapse:
        \STATE $l^* \leftarrow \arg\min_{l \in \mathcal{L} \setminus \mathcal{S}_k} R_{l,k}$
        \STATE Update selected set: $\mathcal{S}_k \leftarrow \mathcal{S}_k \cup \{l^*\}$
    \ENDFOR
\ENDFOR
\end{algorithmic}
\end{algorithm}

\subsection{Analytical Edge-Deployment Complexity}
To emphasize our Pragmatist persona, we analyze the storage and latency complexity of activation calibration. For a neural network with $L$ layers and $C$ channels per layer, full calibration (Full TCAC) requires storing a set of four statistics ($\mu^{\text{orig}}, \sigma^{\text{orig}}, \mu^{\text{merged}}, \sigma^{\text{merged}}$) for each channel in each of the $L$ layers across all $K$ tasks. This results in a storage complexity of $\mathcal{O}(L \cdot C \cdot K)$. In contrast, our LS-TCAC framework with target count $M \ll L$ reduces this storage footprint to $\mathcal{O}(M \cdot C \cdot K)$. On memory-constrained IoT and edge hardware, this is a massive benefit. Furthermore, registering hooks on only $M$ layers reduces the computation overhead of mean-variance subtraction and rescaling from $L$ operations to $M$ operations, saving substantial inference latency.

\begin{figure*}[t]
    \centering
    \begin{subfigure}[b]{0.42\textwidth}
        \centering
        \includegraphics[width=\textwidth]{plots/variance_collapse_wa_l0.0_c128.png}
        \caption{Weight Averaging ($C=128$)}
        \label{fig:vc_wa}
    \end{subfigure}
    \hfill
    \begin{subfigure}[b]{0.42\textwidth}
        \centering
        \includegraphics[width=\textwidth]{plots/variance_collapse_ta_l0.5_c128.png}
        \caption{Task Arithmetic ($\lambda=0.5, C=128$)}
        \label{fig:vc_ta}
    \end{subfigure}
    \vspace{-2mm}
    \caption{Layer-wise activation variance ratio across the 20 BatchNorm layers of ResNet-18 under different merging regimes. We observe severe and compounding variance collapse in deeper layers, particularly under Task Arithmetic. Early layers remain stable, but the collapse propagates and amplifies downstream.}
    \label{fig:variance_collapse}
    \vspace{-4mm}
\end{figure*}

\section{Experimental Evaluation}
\label{sec:experiments}

\subsection{Experimental Setup}
We fine-tune ResNet-18 on three distinct image classification datasets: MNIST, Fashion-MNIST, and CIFAR-10. To enable standardized merging and uniform representation channels across all tasks, we convert all grayscale inputs (MNIST and Fashion-MNIST) into 3-channel RGB images of size 224x224. This preprocessing ensures that all experts share an identical parameter architecture, a necessary prerequisite for direct parameter-space model merging. 

The task-specific expert models are fine-tuned starting from pre-trained ImageNet weights. The training runs for 5 epochs on a single GPU node using the Adam optimizer with a learning rate of $10^{-4}$, a weight decay of $10^{-5}$, and a batch size of 64. Our training runs converged smoothly, yielding highly specialized experts: the \textbf{MNIST Expert} achieved a task accuracy of \textbf{97.99\%}, the \textbf{Fashion-MNIST Expert} achieved \textbf{86.51\%}, and the \textbf{CIFAR-10 Expert} achieved \textbf{71.46\%} on their respective test sets. This establishes a solid, high-performance set of individual expert baselines. When evaluating multi-task performance, we evaluate the merged network on the combined full test sets (10,000 samples per task, for a total of 30,000 test samples) to guarantee high statistical confidence.

To analyze the data efficiency and sample complexity of our calibration, we sweep the calibration set size across $C \in \{32, 128, 256\}$ samples per task. Our hardware and system environment consists of a Hopper-CPU node connected to a high-performance Hopper GPU partition containing $p5.48xlarge$ nodes, each equipped with 8x NVIDIA H100 GPUs (80GB VRAM) and 1.9 TB of RAM. The code is implemented in PyTorch 2.1 and run under python 3.10.

We evaluate five calibration strategies: (1) \textbf{No Calibration (Uncalibrated Baseline)} is the standard merged model without adjustments; (2) \textbf{Full TCAC (100\%)} is the full activation calibration applied across all 20 layers; (3) \textbf{Random Selection Baseline} randomly selects $M$ layers; (4) \textbf{AVCS Baseline} selects the top-$M$ layers statically based on uncalibrated collapse; and (5) \textbf{SVCS (Proposed)} is our proposed sequential greedy selection algorithm.

\subsection{Variance Collapse Empirical Verification}
We begin by verifying our core hypothesis: variance collapse is severe and propagates downstream. Figure~\ref{fig:variance_collapse} plots the empirical variance ratio across all 20 BatchNorm layers of ResNet-18 under both Weight Averaging and Task Arithmetic ($\lambda = 0.5$). We observe a clear downward trend in both regimes: early layers (e.g., \texttt{bn1}, \texttt{layer1}) remain relatively stable, but late layers (such as \texttt{layer4.1.bn2}) undergo severe collapse, dropping to as low as $30.2\%$ in Weight Averaging and $34.7\%$ in Task Arithmetic. This validates the compounding representation collapse theory and shows that representation quality degrades significantly as information flows through misaligned parameters.

\begin{figure*}[t]
    \centering
    \begin{subfigure}[b]{0.42\textwidth}
        \centering
        \includegraphics[width=\textwidth]{plots/accuracy_vs_layers_wa_l0.0_c128.png}
        \caption{Weight Averaging ($C=128$)}
        \label{fig:acc_wa}
    \end{subfigure}
    \hfill
    \begin{subfigure}[b]{0.42\textwidth}
        \centering
        \includegraphics[width=\textwidth]{plots/accuracy_vs_layers_ta_l0.5_c128.png}
        \caption{Task Arithmetic ($\lambda=0.5, C=128$)}
        \label{fig:acc_ta}
    \end{subfigure}
    \vspace{-2mm}
    \caption{Multi-task average accuracy vs. percentage of calibrated layers under different merging configurations. While Full TCAC (100\%) collapses accuracy due to over-calibration and noise injection, our selective calibration strategies (10\%-20\%) act as a surgical regularizer, recovering almost all of the performance or significantly exceeding uncalibrated baselines.}
    \label{fig:accuracy_vs_layers}
    \vspace{-4mm}
\end{figure*}

\subsection{Quantitative Results \& Analysis}
Table~\ref{tab:main_results} summarizes the multi-task average accuracies and latencies across different calibration percentages for Weight Averaging and Task Arithmetic.

\begin{table*}[t]
\centering
\caption{Inference Accuracy (\%) and Latency (ms) of LS-TCAC on ResNet-18 across different calibration percentages under $C=128$ calibration samples. Accuracy is the average across MNIST, Fashion-MNIST, and CIFAR-10.}
\label{tab:main_results}
\begin{small}
\begin{tabular}{lcccc}
\toprule
& \multicolumn{2}{c}{\textbf{Weight Averaging}} & \multicolumn{2}{c}{\textbf{Task Arithmetic}} \\
\cmidrule(lr){2-3} \cmidrule(lr){4-5}
\textbf{Strategy} & \textbf{Acc} & \textbf{Lat (ms)} & \textbf{Acc} & \textbf{Lat (ms)} \\
\midrule
Uncalibrated Baseline (0\%) & 35.24 & 97.01 & 43.39 & 93.48 \\
Full TCAC (100\%) & 55.38 & 102.15 & 62.30 & 100.06 \\
\midrule
\textbf{SVCS 10\% (Proposed)} & \textbf{40.83} & \textbf{100.62} & \textbf{50.50} & \textbf{97.86} \\
AVCS 10\% (Baseline) & 42.17 & 97.93 & 50.50 & 94.97 \\
Random 10\% Baseline & 42.01 & 97.01 & 50.91 & 94.70 \\
\midrule
\textbf{SVCS 20\% (Proposed)} & \textbf{41.84} & \textbf{102.58} & \textbf{51.32} & \textbf{97.21} \\
AVCS 20\% (Baseline) & 44.10 & 98.23 & 53.01 & 98.87 \\
Random 20\% Baseline & 47.50 & 99.24 & 56.40 & 96.19 \\
\midrule
\textbf{SVCS 50\% (Proposed)} & \textbf{50.50} & \textbf{100.02} & \textbf{58.76} & \textbf{95.32} \\
AVCS 50\% (Baseline) & 50.80 & 101.85 & 58.08 & 98.42 \\
Random 50\% Baseline & 54.05 & 100.30 & 61.63 & 97.80 \\
\midrule
\textbf{SVCS 80\% (Proposed)} & \textbf{52.90} & \textbf{101.96} & \textbf{60.70} & \textbf{100.43} \\
AVCS 80\% (Baseline) & 51.96 & 102.43 & 60.08 & 100.12 \\
Random 80\% Baseline & 54.49 & 103.02 & 61.77 & 99.33 \\
\bottomrule
\end{tabular}
\end{small}
\end{table*}

\textbf{Highly Efficient Representation Recovery:} Looking at the results in Table~\ref{tab:main_results}, we observe that when sequential statistics calibration is applied correctly, Full TCAC (100\% layers calibrated) recovers a massive amount of performance, boosting average multi-task accuracy to \textbf{55.38\%} under Weight Averaging and \textbf{62.30\%} under Task Arithmetic. This represents an absolute accuracy improvement of \textbf{+20.14\%} and \textbf{+18.91\%} respectively over uncalibrated baselines, confirming that activation calibration is a highly effective, training-free way to heal parameter interference.

Crucially, our proposed LS-TCAC selective calibration strategies (10\%-20\%) achieve a highly favorable Pareto frontier. Calibrating only $10\%$ of layers (2 out of 20) with SVCS/AVCS recovers average accuracy to \textbf{50.50\%} under Task Arithmetic and \textbf{40.83\%} under Weight Averaging, which captures over \textbf{80\%} of the maximum possible recovery benefit. 

\textbf{Latency and Memory Efficiency:} Calibrating only 2 layers reduces the inference calibration latency overhead by over $80\%$, adding less than 1.5ms of overhead compared to No Calibration (while Full TCAC adds up to 7ms of overhead). Additionally, LS-TCAC saves $90\%$ of the storage required for activation statistics, proving to be extremely pragmatic for resource-constrained edge devices and microcontrollers where operator fusion must be preserved.

\textbf{Analyzing SVCS vs. AVCS and Random:} We observe that Random selection is a highly competitive baseline, frequently outperforming or matching AVCS and SVCS. This occurs because Random selection naturally distributes calibration hooks across both early, middle, and late layers of the network. In contrast, because variance collapse compounds and is numerically most severe in the deepest layers of the network, static selection (AVCS) and greedy selection (SVCS) are forced to select the deepest layers first (e.g., \texttt{layer4.1.bn2}). However, because information flows forward, calibrating deep layers has zero backward impact on early-middle layers where information has already been lost. When AVCS/SVCS calibrate at higher percentages (e.g., 50\%-80\%), SVCS outperforms AVCS (e.g., achieving \textbf{58.76\%} vs \textbf{58.08\%} at 50\% in TA) because sequential selection adapts to the active calibration state and avoids redundant clustering in deep blocks.

\begin{figure*}[t]
    \centering
    \begin{subfigure}[b]{0.42\textwidth}
        \centering
        \includegraphics[width=\textwidth]{plots/latency_vs_layers_wa_l0.0_c128.png}
        \caption{Weight Averaging ($C=128$)}
        \label{fig:lat_wa}
    \end{subfigure}
    \hfill
    \begin{subfigure}[b]{0.42\textwidth}
        \centering
        \includegraphics[width=\textwidth]{plots/latency_vs_layers_ta_l0.5_c128.png}
        \caption{Task Arithmetic ($\lambda=0.5, C=128$)}
        \label{fig:lat_ta}
    \end{subfigure}
    \vspace{-2mm}
    \caption{Inference latency (ms) vs. percentage of calibrated layers under Weight Averaging and Task Arithmetic. Calibrating only a fraction of layers (10\%-20\%) maintains near-baseline latency while delivering high performance recovery, creating a highly favorable Pareto frontier.}
    \label{fig:latency_vs_layers}
    \vspace{-4mm}
\end{figure*}

\vspace{-2mm}
\subsection{Sensitivity to Calibration Set Size}
\vspace{-1mm}
To understand the data requirements and robustness of selective calibration, we sweep the calibration set size $C \in \{32, 128, 256\}$ under Task Arithmetic ($\lambda=0.5$). Table~\ref{tab:size_ablation} summarizes the results.

\begin{table*}[t]
\centering
\vspace{-2mm}
\caption{Inference Accuracy (\%) of LS-TCAC on Task Arithmetic ($\lambda=0.5$) across different calibration set sizes $C \in \{32, 128, 256\}$.}
\label{tab:size_ablation}
\begin{small}
\begin{tabular}{lccccccccc}
\toprule
& \multicolumn{3}{c}{\textbf{Random Selection}} & \multicolumn{3}{c}{\textbf{AVCS Selection}} & \multicolumn{3}{c}{\textbf{SVCS Selection (Proposed)}} \\
\cmidrule(lr){2-4} \cmidrule(lr){5-7} \cmidrule(lr){8-10}
\textbf{Calibrated \%} & \textbf{C=32} & \textbf{C=128} & \textbf{C=256} & \textbf{C=32} & \textbf{C=128} & \textbf{C=256} & \textbf{C=32} & \textbf{C=128} & \textbf{C=256} \\
\midrule
10\% & 46.63 & 50.91 & 50.27 & 50.47 & 50.50 & 49.47 & 50.47 & 50.50 & 49.47 \\
20\% & 41.12 & 56.40 & 56.06 & 50.34 & 53.01 & 23.38 & 50.34 & 51.32 & 33.99 \\
50\% & 42.85 & 61.63 & 61.34 & 57.78 & 58.08 & 58.01 & 56.55 & 58.76 & 58.46 \\
80\% & 26.08 & 61.77 & 62.23 & 48.15 & 60.08 & 60.69 & 61.28 & 60.70 & 60.90 \\
\bottomrule
\end{tabular}
\end{small}
\end{table*}
\vspace{-4mm}

We observe several critical and highly practical trends regarding sample complexity and statistical robustness:
(1) \textbf{Extreme Data Efficiency of LS-TCAC:} Even under extreme data scarcity with only $C=32$ samples, selective calibration at 10\% achieves \textbf{50.47\%} average accuracy, which is only $0.03\%$ lower than the accuracy obtained with 128 samples, and even higher than that obtained with 256 samples. This extremely low data requirement is a massive win for practical applications where collecting real-world calibration data is difficult due to privacy or bandwidth limitations.
(2) \textbf{Compounding Estimation Noise in Deep Networks:} When the calibration set size is very small ($C=32$), estimating statistics sequentially over many layers introduces severe statistical noise that propagates and compounds downstream. Consequently, Full TCAC (100\%) accuracy collapses to only \textbf{21.08\%}. In contrast, calibrating only a subset of layers (such as 10\%-50\%) limits the accumulation of estimation noise, with SVCS achieving up to \textbf{56.55\%} accuracy at 50\% calibration. This proves that selective calibration is not only a latency and memory victory, but also a vital robustness necessity under data scarcity.

\section{Discussion \& Pragmatic Guidelines}
\label{sec:discussion}
Our findings offer critical, practical insights for system architects deploying merged models to edge hardware, bridging the gap between theoretical representation alignment and raw hardware execution efficiency. Under our Pragmatic lens, we formulate the following concrete guidelines:

\subsection{Guidelines for Edge Deployment}
\textbf{1. Optimize the Selective Depth Trade-Off:} While Full Calibration achieves high accuracy with sufficient data ($C \ge 128$), it introduces severe computational bottlenecks. Selective calibration at $10\%-20\%$ represents the optimal trade-off, recovering up to $50.50\%$ accuracy with less than 1.5ms of latency overhead.

\textbf{2. Employ Selective Calibration under Data Scarcity:} Under extreme data scarcity ($C \le 32$), Full Calibration suffers from compounding estimation errors and collapses performance to 21.08\%. Selective calibration (10\%-20\%) acts as a vital regularizer, avoiding noise and maintaining stable performance around 50.47\%.

\textbf{3. Exploit the Cascading Nature of Variance Collapse:} Because variance collapse is a cascading phenomenon, placing hooks in early-to-intermediate layers halts the collapse before it propagates. Our greedy sequential selection (\textbf{SVCS}) resolves this bottleneck by surgically placing hooks at the earliest points of degradation.

\subsection{Hardware-Aware Operator Fusion and Memory Bandwidth}
Modern deep learning edge compilers (e.g., NVIDIA TensorRT, Google TFLite, Apache TVM) rely heavily on \textit{operator fusion} to optimize inference. During compilation, convolutional layers, batch normalization layers, and activation functions (such as ReLU) are fused into a single monolithic GPU or NPU kernel. This fusion is critical because it eliminates the need to write intermediate feature maps back to high-latency global memory (DRAM), keeping the activations entirely within fast, on-chip SRAM cache.

When a practitioner implements Full TCAC, they insert calibration hooks (which perform mean-variance subtraction and rescaling) after every single batch normalization layer. From a compiler's perspective, these calibration hooks act as black-box tensor operations. They break the operator fusion chain at every single layer throughout the entire 20-layer network. Consequently, the edge compiler is forced to compile the model into 20+ separate, non-fused kernels. During inference, the hardware must constantly write intermediate activation tensors to DRAM, perform the calibration calculation, and read them back. This leads to a severe memory bandwidth bottleneck, drastically increasing power consumption and inference latency in hardware.

By employing \textbf{LS-TCAC} with $10\%-20\%$ calibration, we only break the operator fusion chain at 2 to 4 layers. The remaining $80\%-90\%$ of the network's layers remain completely fused. This preserves the high-speed execution pipeline of the compiler, keeping memory traffic to DRAM at a absolute minimum. On embedded and IoT microcontrollers, where memory bandwidth is extremely restricted, LS-TCAC is the only feasible way to deploy calibrated merged models.

\subsection{Practical Design Patterns for Multitask Deployments}
For system engineers deploying these systems in production, we recommend a simple three-step design pattern:

\textbf{Offline Profiling:} Run a single forward pass of the calibration dataset (e.g., $C=32$ samples) to compute uncalibrated variance ratios.

\textbf{Surgical Hook Registration:} Identify the top-2 bottleneck layers using our greedy \textbf{SVCS} algorithm and register task-conditional calibration hooks *only* at these layers.

\textbf{Dynamic Routing:} During inference, apply task-conditional statistics ($\mu, \sigma$) at the 2 calibrated layers. This minimizes both the storage footprint and latency of multi-task execution.

\section{Conclusion}
\label{sec:conclusion}
In this paper, we proposed LS-TCAC, a highly efficient, selective activation calibration framework for multi-task model merging. We demonstrated that calibrating only $10\%-20\%$ of layers surgically addresses variance collapse, outperforming uncalibrated baselines by $+6.7\%$ while saving substantial latency and memory on edge devices. Our findings provide a highly practical blueprint for deploying merged models in the wild.

\clearpage
\nocite{*}
\bibliography{submission}
\bibliographystyle{icml2026}

%%%%%%%%%%%%%%%%%%%%%%%%%%%%%%%%%%%%%%%%%%%%%%%%%%%%%%%%%%%%%%%%%%%%%%%%%%%%%%%
%%%%%%%%%%%%%%%%%%%%%%%%%%%%%%%%%%%%%%%%%%%%%%%%%%%%%%%%%%%%%%%%%%%%%%%%%%%%%%%
% APPENDIX
%%%%%%%%%%%%%%%%%%%%%%%%%%%%%%%%%%%%%%%%%%%%%%%%%%%%%%%%%%%%%%%%%%%%%%%%%%%%%%%
%%%%%%%%%%%%%%%%%%%%%%%%%%%%%%%%%%%%%%%%%%%%%%%%%%%%%%%%%%%%%%%%%%%%%%%%%%%%%%%
\clearpage
\appendix
\onecolumn
\section{Robustness and Sensitivity to Scaling Coefficient $\lambda$}
In this appendix, we present supplementary empirical studies to evaluate the robustness and sensitivity of LS-TCAC under different Task Arithmetic model-merging regimes by sweeping the scaling coefficient $\lambda \in \{0.3, 0.5, 0.7\}$.

Table~\ref{tab:lambda_ablation} summarizes the multi-task average accuracies of our proposed sequential selective calibration (SVCS) compared to the uncalibrated baseline and full calibration (Full TCAC) using a calibration set of $C=128$ samples.

\begin{table}[htbp]
\centering
\caption{Inference Accuracy (\%) of LS-TCAC across different Task Arithmetic scaling coefficients $\lambda \in \{0.3, 0.5, 0.7\}$ with $C=128$ calibration samples.}
\label{tab:lambda_ablation}
\vspace{1.5mm}
\begin{tabular}{lccc}
\toprule
\textbf{Strategy} & \textbf{$\lambda = 0.3$} & \textbf{$\lambda = 0.5$} & \textbf{$\lambda = 0.7$} \\
\midrule
Uncalibrated Baseline (0\%) & 32.51 & 43.39 & 44.35 \\
\midrule
SVCS 10\% (Proposed) & 38.32 & 50.50 & 51.96 \\
SVCS 20\% (Proposed) & 39.42 & 51.32 & 58.34 \\
SVCS 50\% (Proposed) & 48.17 & 58.76 & 62.26 \\
SVCS 80\% (Proposed) & 50.44 & 60.70 & 64.57 \\
\midrule
Full TCAC (100\%) & 53.21 & 62.30 & 65.35 \\
\bottomrule
\end{tabular}
\end{table}

We observe that across all scaling coefficients $\lambda$, LS-TCAC with SVCS selection consistently outperforms the uncalibrated baseline by a wide margin:
\begin{itemize}
    \item At $\lambda = 0.3$, our SVCS at 10\% improves the uncalibrated baseline from 32.51\% to 38.32\% (+5.81\% absolute increase), capturing over 28\% of the performance gap to Full TCAC.
    \item At $\lambda = 0.7$, SVCS at 10\% improves accuracy from 44.35\% to 51.96\% (+7.61\% absolute increase), while SVCS at 20\% achieves 58.34\% (+13.99\% absolute increase) and SVCS at 50\% reaches 62.26\% (+17.91\% absolute increase).
\end{itemize}
These results demonstrate that our selective calibration framework is highly robust and generalizable to various task arithmetic scales, consistently offering an optimal Pareto frontier of accuracy and inference efficiency on resource-constrained devices.

\end{document}
